\section{Related Work}

StructFlow connects two lineages of flow-based modeling: causal autoregressive flows for structural modeling, and conditional flow matching for Bayesian inference.

\paragraph{Causal Generative Modeling.}
\citet{khemakhem2021causal} formalized the intersection of autoregressive flows and structural causal models (SCMs) with CAREFL: autoregressive architectures are structurally equivalent to SCMs, so interventional and counterfactual queries can be answered by fixing the base distribution (exogenous noise). Our \textbf{Causal Alignment} mode adopts this structural perspective but implements the transformation with continuous flow matching, retaining the causal interpretability while gaining the flexibility of ODE-based generation; related structural constraints appear in \citet{wehenkel2021graphical} and \citet{pawlowski2020deep}. A parallel line answers the same queries with other generative primitives: CausalGAN \citep{kocaoglu2018causalgan} arranges adversarial generators along the causal graph, neural causal models \citep{xia2021causal} establish expressiveness and identifiability guarantees for graph-structured generators, and Modular-DCM \citep{rahman2024modular} scales their training by reusing pretrained components. Closest to StructFlow are the diffusion-based causal models Diff-SCM \citep{sanchez2022diffusion} and DCM \citep{chao2023interventional}, which train one diffusion model per node, conditioned on its parents, and answer counterfactuals by recovering the exogenous noise through deterministic DDIM inversion---structurally, our causal-alignment mode with a diffusion backbone. All of these cover the causal direction only: no density evaluation, no amortized model of the inference direction (posterior queries require expensive per-query procedures), and no treatment of the denser dependencies explaining-away induces. StructFlow complements them by compiling the same graph specification into either direction.

\paragraph{On backbones.}
\label{par:backbone}
Flow matching and diffusion are the same object class up to reparameterization: every diffusion model admits a probability-flow ODE with the same marginals, DDIM is a discretization of that ODE, and linear-interpolant flow matching is a time-and-scale reparameterization of a variance-preserving diffusion \citep{song2021scorebased, song2021ddim, gao2024diffusionmeetsflow}. Our contribution---the compiler and the block-triangular architecture over an arbitrary DAG---is therefore orthogonal to the backbone. The one place the backbone does bite is abduction: DDIM inversion carries a first-order step-size bias where a CNF's inversion error is solver tolerance at fixed weights (Proposition~\ref{prop:abduction-error}). We test both claims directly by swapping a diffusion backbone into our compiler (Section~\ref{sec:backbone-ablation}).

\paragraph{Flow Matching and Amortized Inference.}
\citet{wildberger_flow_2023} learn posteriors $p(\theta|x)$ for scientific simulators with flow matching, but in a black-box manner. Most relevant is \citet{jeong_conditional_2025}, who introduced \textbf{dynamic conditioning}: the velocity field is conditioned on time-dependent trajectories of the observed variables rather than static inputs, yielding a triangular transport map respecting the posterior's dependency structure. We generalize their method from a single observation-to-target split to arbitrary PGMs.
Among other structured approaches, \citet{weilbach2020structured} construct structured CNFs for amortized inference via a minimally faithful inversion of the graph, but target the inference direction only; \citet{javaloy2023causal} develop causal normalizing flows with identifiability guarantees but focus on the generative direction; \citet{chen2023structured} use masked structured networks for density and causal-effect estimation without amortized posterior sampling; and \citet{baptista2024representation} study monotone triangular transport maps whose sparsity pattern reflects the conditional-independence structure of the target. Each covers a \emph{subset} of the capabilities we target (Table~\ref{tab:capability}, Appendix~\ref{app:additional-experiments}). We do not claim a new flow primitive; our contribution is a single automated compiler that unifies the causal and inference directions for \emph{arbitrary} DAGs---including the moralization needed to handle explaining-away at colliders, and both static and dynamic conditioning paradigms.
